% Necessarily a separate file because babel interferes with amsmath. The pdfs will be concatenated manually.

\documentclass[11pt]{article}

\title{קבוצת ז׳וליה של $z^2+1/4$ היא כמעט-קמורה}

\date{Month and Year}

\usepackage[utf8]{inputenc}
\usepackage{csquotes}
\usepackage[hebrew,english]{babel}
\usepackage{biblatex}
\usepackage{amsmath}
\usepackage{tikz}
\usepackage{indentfirst}
\usepackage{lipsum}  

\parskip 2mm
\addtolength{\textwidth}{1.4in} \addtolength{\oddsidemargin}{-0.5in}
\addtolength{\evensidemargin}{-0.8in}
\addtolength{\topmargin}{-0.5in} \addtolength{\textheight}{0.1in}

\DeclareUnicodeCharacter{05D0}{\hebalef}
\DeclareUnicodeCharacter{05D1}{\hebbet}
\DeclareUnicodeCharacter{05D2}{\hebgimel}
\DeclareUnicodeCharacter{05D3}{\hebdalet}
\DeclareUnicodeCharacter{05D4}{\hebhe}
\DeclareUnicodeCharacter{05D5}{\hebvav}
\DeclareUnicodeCharacter{05D6}{\hebzayin}
\DeclareUnicodeCharacter{05D7}{\hebhet}
\DeclareUnicodeCharacter{05D8}{\hebtet}
\DeclareUnicodeCharacter{05D9}{\hebyod}
\DeclareUnicodeCharacter{05DA}{\hebfinalkaf}
\DeclareUnicodeCharacter{05DB}{\hebkaf}
\DeclareUnicodeCharacter{05DC}{\heblamed}
\DeclareUnicodeCharacter{05DD}{\hebfinalmem}
\DeclareUnicodeCharacter{05DE}{\hebmem}
\DeclareUnicodeCharacter{05DF}{\hebfinalnun}
\DeclareUnicodeCharacter{05E0}{\hebnun}
\DeclareUnicodeCharacter{05E1}{\hebsamekh}
\DeclareUnicodeCharacter{05E2}{\hebayin}
\DeclareUnicodeCharacter{05E3}{\hebfinalpe}
\DeclareUnicodeCharacter{05E4}{\hebpe}
\DeclareUnicodeCharacter{05E5}{\hebfinaltsadi}
\DeclareUnicodeCharacter{05E6}{\hebtsadi}
\DeclareUnicodeCharacter{05E7}{\hebqof}
\DeclareUnicodeCharacter{05E8}{\hebresh}
\DeclareUnicodeCharacter{05E9}{\hebshin}
\DeclareUnicodeCharacter{05EA}{\hebtav}
\usepackage{amsfonts}

\begin{document}

\selectlanguage{hebrew}
\pagenumbering{gobble}
\begin{abstract}
{\large

    לכל מספר מרוכב
    $z \in \mathbb C$
    ,
    נגדיר באופן אינדוקטיבי סדרה
$a_n=a_n(z)$
 על ידי
 כלל הנסיגה
    $$a_{n+1} = a_n ^2 + 1/4,$$
    עם נקודת ההתחלה
        $$a_0 = z.$$
    
    \textbf{אגן המשיכה לאינסוף}
    של הפולינום
    $$f(z)=z^2+1/4$$        
    % תהי $A \subset \mathbb C$
    הוא
    קבוצת נקודות ההתחלה 
    $z$
    שעבורן
    $a_n(z) \to \infty$,
    כלומר
הקבוצה
     $$A_\infty(f)=\{z\in \mathbb C: f^{\circ n}(z)=((z^2+1/4)^2+1/4)^2+\dots \to \infty\}.$$
אגן המשיכה של 
$f$
הוא
תחום ששפתו 
$\mathcal J=\partial A_\infty(f)$
היא עקומת ז'ורדן.

העקומה $\mathcal J$
היא פרקטל:
ממד האוסדורף של $\mathcal J$
הוא בקירוב 
$1.0812$.
% https://abel.math.harvard.edu/~ctm/papers/home/text/papers/dimIII/dimIII.pdf appendix A


עבודה זו מוכיחה שאגן המשיכה של 
$f$
 הוא 
\textbf{כמעט-קמור}:
ניתן לחבר כל שתי נקודות בו באמצעות עקומה
$\eta$
שמוכלת באגן המשיכה, 
ושהאורך שלה גדול רק פי קבוע מאורך הקו הישר שמחבר בין הנקודות.

    
    % \textbf{אגן המשיכה לאינסוף} של פולינום $f$ הוא אוסף המספרים המרוכבים ששואפים לאינסוף תחת הפעלה נשנית של 
    % $f$.
    
    %  עבודה זו עוסקת באגן המשיכה לאינסוף של הפולינום 
    %  $f(z)=z^2+1/4$,
    %  כלומר בקבוצה 
    %  $$A=\{z\in \mathbb C: f^{\circ n}(z)=((z^2+1/4)^2+1/4)^2+\dots \to \infty\}.$$
     
    }
    
\end{abstract}

\selectlanguage{english}
\begin{titlepage}
    \begin{center}
        \vspace*{0.05cm}
        \includegraphics[width=0.8\textwidth]{figures/Logo_Exact_Sciences_Hebrew_Black.jpeg}\\
        \selectlanguage{hebrew}
        \vspace{3cm}
        % {\huge קבוצת ז'וליה של $z^2+1/4$ היא כמעט-קמורה \par}
        {\huge אגן המשיכה לאינסוף של $z^2+1/4$ הוא כמעט-קמור \par}
        \vspace{1cm}
        {\large חיבור זה הוגש כחלק מהדרישות לקבלת התואר\\
        \textquotedblright מוסמך אוניברסיטה\textquotedblleft\text{ }\textendash \text{ }.cS.M}\\
        \vspace{2cm}
        { על ידי}\\
        \vspace{0.3cm}
        {\large {עמנואל סגל}}\\
        \vspace{1.5cm}
        {\large העבודה נכתבה בהדרכתם של}\\
        \vspace{0.3cm}
        {\huge דר' אולג עברי} \\
        {\huge פרופ' מיכאיל סודין\\
        \vspace{1cm}
        {\Large בית הספר למדעי המתמטיקה \\ אוניברסיטת תל אביב}}
    \end{center}
\end{titlepage}

\end{document}